\section{Pruebas}
\begin{frame}{Muestra el resultado y explícalo}
\begin{columns}[c,onlytextwidth]
\column{0.54\textwidth}
\centering
\begin{tabular}{@{}lrr@{}}
\toprule
& Habitual & Revisado\\
\midrule
Tiempo (min) & 24 & 18\\
Errores (\%) & 8 & 5\\
Tareas & 40 & 40\\
\bottomrule
\end{tabular}
\par\vspace{3mm}{\scriptsize\color{muted}Valores ilustrativos; usa tus propios datos.}
\column{0.38\textwidth}
{\Huge\bfseries\color{accent}25\%}\par\vspace{2mm}
menos tiempo en este ejemplo.\par\vspace{4mm}
Explica si la diferencia es relevante y cuál es la incertidumbre de la estimación.
\end{columns}
\end{frame}

\begin{frame}{Distingue el hallazgo de sus límites}
\begin{columns}[T,onlytextwidth]
\column{0.46\textwidth}
\begin{block}{Lo que respaldan las pruebas}
Expón el hallazgo principal en una frase. Conéctalo con la pregunta de
investigación y las pruebas de la diapositiva anterior.
\end{block}
\column{0.46\textwidth}
\textbf{Lo que sigue sin resolverse}\par\vspace{3mm}
\begin{itemize}\setlength{\itemsep}{3mm}
\item Alcance y tamaño de la muestra.
\item Supuestos de medida.
\item Generalización a otros contextos.
\end{itemize}
\end{columns}
\end{frame}
